\subsection{Structural definition and role of the Barbero--Immirzi parameter}
\label{sec:ii-definicion-barbero}

\begin{definition}[Oriented evaluation of hexadic--octadic incidence]
The parameter \(\gamma_{\HMT}\) is the conjugate evaluation of the
fundamental chain \(c_{12}^{+}\), completed by its angular
component. On the domain \(0<q_\pm<1\), its definition
\eqref{eq:gamma} takes the form
\begin{equation}
 \gamma_{\HMT}
 =\frac{\pi AC^*}{180}
  +[\mathscr R(c_{12}^{+})](q_-)
  -[\mathscr R(c_{12}^{+})](q_+).
 \label{eq:ii-definicion-barbero-evaluacion}
\end{equation}
The marked incidence supplies the event spaces; the oriented
chain fixes their multiplicities; the channels in
\eqref{eq:ii-canales} perform the evaluation.
\end{definition}

The determination jointly uses the cardinalities \(90,120\),
the twelve sector contributions, and the return seam. The
pole coefficient \(1/24\) and Diophantine minimality verify
the compatibility of these weights with the normalization. Their provenance
is preserved in \(c_{12}^{+}\), before the channels are evaluated. Angular
inversion exchanges \(q_+\) and \(q_-\) and produces
\(\gamma_{\HMT}(A,-C^*)=-\gamma_{\HMT}(A,C^*)\): the dimensionless
value transports an orientation of the incidence.

In the geometric realization, this same coefficient enters
\eqref{eq:holst} and the area increment
\(\gamma_{\HMT}a_0\ell_P^2\). The Holst action uses
\(\gamma_{\HMT}\ne0\); the positive areal magnitude uses
\(|\gamma_{\HMT}|\). Proposition~\ref{prop:ii-costura-barbero}
makes the coordinated change of orientation explicit. Thus, the structure
of the parameter explains what it weights in the incidence, what it preserves
under transport, and how it participates in the torsional response and in
the geometric resolution of information.
